The main result is architectural rather than mechanistic. Under the matched-output evaluation, a coding-agent-based system organized around repeated choices among planning, drafting, and reviewing is preferred to one-shot generation, a prescribed writing pipeline, and adaptive task decomposition. The independent product metrics are more mixed: \condAgenticCogWriter\ is strongest on two benchmarks but not the third. Together, these results support process-organized execution as a useful way to structure an agent that must produce a long document, without implying that it is uniformly superior under every quality measure.

The interventions narrow what can explain that result. The explicit goal network is a natural feature to expect to matter, especially given its motivation from writing theory, yet removing it does not significantly worsen pairwise preference under the pre-specified tests. Flexible process order is similarly non-essential in the tested form: fixing planning, drafting, and reviewing to a repeated cycle does not significantly worsen preference, and the unrestricted host agent chooses that cycle in most completed runs anyway. The trace split also gives no indication that departures from the cycle are where the performance difference comes from. These findings make it difficult to attribute the result to either explicit goal bookkeeping or flexible scheduling alone.

The execution variants also caution against a simple ``multi-agent is better'' explanation. \condSingleWriter\ performs substantially worse, but that variant changes delegation, role context, control granularity, and output length together. The narrower \condSingleContext\ comparison preserves the process loop and does not show a significant pooled difference in its exploratory replication, while still differing in gateway and transport. Length- and compute-matched analyses preserve important directions but condition on outcomes of execution. The remaining candidates are therefore properties shared across the successful process-organized variants, such as decomposing writing into explicit activities, persisting state between those activities, and repeatedly returning control to a common host agent. Direct interventions on these shared properties are needed before treating any of them as the causal mechanism.

For agent design, the broader implication is that long-form writing may not require a bespoke end-to-end writing model so much as an appropriate control architecture on top of a capable general-purpose harness. Coding-agent-style systems already provide persistent files, tool use, and subagent execution; \condAgenticCogWriter\ uses those generic capabilities to turn writing into a stateful loop rather than a single generation request. This perspective is relevant beyond standalone writing benchmarks. In autonomous research systems, for example, report or paper generation follows other agentic work such as literature review, coding, and experimentation~\cite{schmidgall2025agentlab,lu2024aiscientist}. A writing architecture that can be embedded into the same harness is therefore a natural systems abstraction.

Cognitive writing theory plays a specific role in this account. It does not validate the architecture by analogy to human cognition, and our traces should not be interpreted as human writing processes. Instead, the theory supplies a vocabulary for separating design choices that agent implementations often bundle together. Turning those choices into interventions is useful precisely because some intuitively important components, such as explicit goal representation and flexible ordering, do not explain the observed result on their own.
